% Source: https://artofproblemsolving.com/community/c3h1351709_agh_i_cant_find_the_sixth_rusczyk_number
% Stable identity: aops-topic-1351709-q01
(a) The remainder of $n$ on division by $y$ is $y-1$ exactly when $n\equiv-1\pmod y$, that is, when $y\mid n+1$. So $n$ has the property for all $y=2,\dots,k$ exactly when $n+1$ is a common multiple of $2,3,\dots,k$, i.e. a multiple of $L_k=\operatorname{lcm}(2,\dots,k)$. The smallest positive such $n$ has $n+1=L_k$ (the smallest positive multiple), so $N_k=L_k-1$.

(b) The least common multiples are $2,6,12,60,60,420,840,2520,2520$ for $k=2,\dots,10$, so
\[
N_2,\dots,N_{10}=1,\ 5,\ 11,\ 59,\ 59,\ 419,\ 839,\ 2519,\ 2519 .
\]
(c) \textbf{No.} $N_2=1$ is not prime, and $N_9=N_{10}=2519=11\cdot229$ is composite. (Also, the numbers $7$ and $119$ that a student listed as terms are not of this form at all: $7$ leaves remainder $1$, not $2$, on division by $3$, and $119=7\cdot17$ is composite.) The conjecture is $\boxed{\text{false}}$.
